% --- Archon physics formalization source begin ---
% archon:physics
% archon:covers PhyXMiniProblems/problem_phyx_mini_0012.lean
% archon:source-report reports/phyx_mini/problem_phyx_mini_0012.source.json

\chapter{Physics problem phyx\_mini\_0012}
\label{ch:PhyXMiniProblems_problem_phyx_mini_0012}

\paragraph{Problem source.}
\#\# Physical scenario

An optical fiber has an index of refraction \textbackslash{}( n \textbackslash{}) and diameter \textbackslash{}( d \textbackslash{}). It is surrounded by vacuum. Light is sent into the fiber along its axis as shown in figure.

\#\# Question

Find the smallest outside radius \textbackslash{}( R\_\{\textbackslash{}text\{min\}\} \textbackslash{}) permitted for a bend in the fiber if no light is to escape.

\#\# Answer choices

- A: A: \textbackslash{}( \textbackslash{}frac\{nd\}\{2n+1\} \textbackslash{})

- B: B: \textbackslash{}( \textbackslash{}frac\{nd\}\{2n-1\} \textbackslash{})

- C: C: \textbackslash{}( \textbackslash{}frac\{nd\}\{n-1\} \textbackslash{})

- D: D: \textbackslash{}( \textbackslash{}frac\{nd\}\{n+1\} \textbackslash{})

\#\# Figure caption (auxiliary; use the image as primary evidence)

The image depicts a section of a curved, hollow pipe or tube with a circular arc shape. 

- The arc is shown as a light blue, semi-transparent band with a uniform inner and outer radius.
- The pipe has an outer radius marked as \textbackslash{}( R \textbackslash{}), and the thickness of the pipe is labeled as \textbackslash{}( d \textbackslash{}).
- Several blue arrows are drawn at one end of the pipe, pointing upwards, indicating a direction of flow or force.
- A black line extends from the center of the curvature to the inner edge of the arc, marked with \textbackslash{}( R \textbackslash{}).
- There is a vertical black line at the left end of the pipe, indicating its thickness, with the letter \textbackslash{}( d \textbackslash{}) next to it.

Overall, the image is a diagrammatic representation, possibly referring to properties of fluid flow or structural analysis through a curved pipe section.

\#\# Dataset metadata

Geometrical Optics; Physical Model Grounding Reasoning, Spatial Relation Reasoning

\paragraph{Recorded answer/context.}
C: C: \textbackslash{}( \textbackslash{}frac\{nd\}\{n-1\} \textbackslash{})

\paragraph{Figure/image path.}
/root/proposal\_for\_physic/science-mango/phyx\_mini\_run/phyx\_data/test\_image/12.png

\paragraph{Formalization target.}
create a compiling Lean file with sorry bodies at `PhyXMiniProblems/problem\_phyx\_mini\_0012.lean`. The Lean declarations must preserve the physical quantities, dimensions or dimensional roles, figure labels, governing-law hypotheses, and final relation expressed by this problem.
Use Mathlib/Physlib names found through LeanExplore where available. If a physics API is missing, introduce faithful local abstractions rather than scalar placeholder aliases.

\begin{theorem}[Physics formalization target]
\label{thm:physics:phyx_mini_0012:target}
\lean{PhyXMini0012.minimum_outside_radius}
\uses{def:physics:phyx-mini-0012:phyxmini0012-fiberlength, def:physics:phyx-mini-0012:phyxmini0012-bentopticalfiber, def:physics:phyx-mini-0012:phyxmini0012-axialrayconfinementmodel}
For a vacuum-clad fiber of index \(n>1\), diameter \(d>0\), and outside
bend radius \(R\), suppose the limiting axial ray has incidence sine
\((R-d)/R\) and is confined exactly when this is at least the critical sine
\(1/n\).  Then the least confining outside radius is
\(R_{\min}=nd/(n-1)\).
\end{theorem}
\begin{proof}
Because \(n>1\) and \(d>0\), the candidate \(R_0=nd/(n-1)\) is larger than
\(d\).  Substitution gives
\((R_0-d)/R_0=1/n\), so the total-internal-reflection equivalence makes
\(R_0\) confining.  If \(R>d\) is any confining radius, the same equivalence
and geometry give \(1/n\leq(R-d)/R\).  Multiplication by the positive
quantities \(n\), \(R\), and \(n-1\) rearranges this to
\(nd/(n-1)\leq R\).  Hence \(R_0\) is the stated least element, and uniqueness
of a least element identifies it with the supplied minimum radius.
\end{proof}

% --- Archon declaration topology begin ---
\section{Declaration topology}
\begin{definition}[Fiber Length]
  \label{def:physics:phyx-mini-0012:phyxmini0012-fiberlength}
  \lean{PhyXMini0012.FiberLength}
  A nonnegative length readout carrying Physlib's physical length dimension
\end{definition}
\begin{proof}
  This is the declared model definition of Fiber Length.
\end{proof}
\begin{definition}[Bent Optical Fiber]
  \label{def:physics:phyx-mini-0012:phyxmini0012-bentopticalfiber}
  \lean{PhyXMini0012.BentOpticalFiber}
  \uses{def:physics:phyx-mini-0012:phyxmini0012-fiberlength}
  The material and geometric data of the vacuum-clad optical fiber
\end{definition}
\begin{proof}
  This is the declared model definition of Bent Optical Fiber.
\end{proof}
\begin{definition}[Axial Ray Confinement Model]
  \label{def:physics:phyx-mini-0012:phyxmini0012-axialrayconfinementmodel}
  \lean{PhyXMini0012.AxialRayConfinementModel}
  \uses{def:physics:phyx-mini-0012:phyxmini0012-fiberlength}
  The least information about the axial family of geometrical-optics rays needed to discuss confinement in a bend. At a proposed outside radius, the limiting ray is the axial ray having the smallest incidence angle at the outer wall
\end{definition}
\begin{proof}
  This is the declared model definition of Axial Ray Confinement Model.
\end{proof}
% --- Archon declaration topology end ---
% --- Archon physics formalization source end ---
